
\subsection{위협 모델 시뮬레이션 및 추가적인 보안 조치}

\begin{figure}[H]
    \centering
    \includegraphics[width=0.8\linewidth]{figure15-1.png}
    \includegraphics[width=0.8\linewidth]{figure15-2.png}
    \caption{Simulated adversary view}
\end{figure}

본 시스템은 소유권 증명을 위해 (증명 - 소유자) 맵핑 정보를 보관한다는 한계를 지니고있다. 이 장에선 데이터베이스가 모두 유출되었다는 가정 하에 공격자 시점에서의 추론 공격(Inference Attack) 시뮬레이션을 수행하였다. 실험에 사용된 OWNERSHIP 테이블을 PROOF와 결합하여 각 DP가 가진 데이터들을 분류 및 집계하였고 속성에 따라 각각의 그래프로 시각화 하였다. Figure 15는 특정 개인(DP)으로 부터 알아낼 수 있는 모든 정보를 시각화 한 결과이다. 공격자는 특정 DP가 소유한 모든 데이터의 종류, 수량과 범위의 분포을 알 수 있으며, 이에 기반해 개인에 대한 종합적인 건강 정보를 추론할 수 있다. 그래프에서 나타나듯, 공격자는 개별 튜플의 정확한 평문 수치를 알 수는 없으나, 붉은 선과 푸른 선 사이의 영역에 데이터 원문이 존재한다는 사실은 확인할 수 있다. 본 시스템에서 사용되는 범위 정보는 실제 의학적 참고 수치에 기반하였기 때문에, 이 정도 정보만으로도 해당 개인의 유병 내역(예: 당뇨, 고혈압 등)을 쉽게 예측할 수 있다. 

\bigskip

Figure 15의 실험결과를 해석하면 다음과 같다: 식별자 '0a3f578b...10505bcb'를 가진 개인(DP)은 30개의 속성에 대해 총 1066개의 데이터를 가지고 있었다. 이 중 공복혈당(Fasting Blood Sugar)은 0과 최대치 사이에 37개의 데이터가 분포해 있음을 알 수 있었다. 이 중 당뇨로 판정되는 125 이상의 데이터는 총 14개를 확인하였다. (참고로 해당 실험에 사용된 데이터셋은 난수에 기반해 생성되어 이같은 무작위적인 분포를 보였지만, 실제 데이터의 경우 훨씬 더 일관된 추세를 보일 것으로 예상된다.) 종합하면 공격자는 이 개인이 어떤 시점에 검진을 받았는지, 어떤 검사를 진행했는지, 그리고 데이터의 범위 분포 및 변화의 추이를 파악 할 수 있다. 나아가 이에 기반해 어떤 진단을 받았는지, 특정 시기의 건강 패턴이나 생활 습관이 어떻게 변화했는지 등을 추론해낼 수 있다. 여기에 전문적인 의학적 지식을 동원한다면 다른 속성 정보를 결합해 해당 개인에 대한 더 정확하고 종합적인 건강 정보를 추론해낼 수 있을 것이다.

\bigskip

이러한 구조적 위협을 완화하고 시스템을 고도화하기 위해, 본 연구는 다음과 같은 추가적인 보안 전략의 적용을 제안한다.

\bigskip

첫째, 더미 데이터 주입을 통한 교란이다. DI 혹은 DP측에서 최초로 증명을 생성할 때, 실제 데이터와 다른 범위에 속하는 데이터를 생성, 이에 대한 더미 증명을 생성한다. 이처럼 모든 증명에 대해, 수학적으로는 참이지만 실존하지는 않는 거짓 양성(False Positive) 증명을 1대 1로 생성하여 SS에 전달한다. 여기서 어느 쪽이 진짜 증명인지는 오직 이를 생성한 DI 혹은 DP만이 알고 있으며, SS 조차도 이를 알 수 없다. 데이터 수요자(DC)는 데이터 구매 결정 후 소유자로부터 데이터 원문을 넘겨받는 시점에서 진짜 증명을 구분할 수 있게 한다. 이 경우 공격자가 소유권 정보를 탈취하여 분석하더라도, 가짜 증명으로 발생한 노이즈로 인해 특정 개인의 병력을 확정적으로 추론하는 것이 불가능해진다. 다만 이러한 구조는 공격자 뿐만 아니라 수요자의 검색 결과에게도 노이즈를 추가하게 되므로 구매 과정에서 이를 효율적으로 필터링 해내는 추가적인 시스템이 요구된다.

\bigskip

둘째, 비공모 다중 서버 기반의 비밀 분산(Secret Sharing)이다. 검색 서버를 서로 공모하지 않는 두 개 이상의 독립된 서버로 분리하고, OWNERSHIP 테이블의 식별자 정보를 안전한 다자간 컴퓨팅(MPC) 프로토콜을 통해서만 연산되도록 분산 저장하는 방식이다. 이는 현실적으로 적용 가능한 가장 검증된 방법이며, 더미 데이터 주입과 달리 데이터 수요자 측에 부담을 주지 않는다는 장점이 있다. 공격자는 특정 임계값을 넘는 숫자의 서버를 해킹하지 않는 이상 복호화 키를 복원할 수 없으며, 이를 통해 단일 서버가 장악되더라도 소유권 정보가 복원되지 않는 강력한 기밀성을 확보할 수 있다. 단, 이는 SS 측의 연산 오버헤드를 크게 증가시키는 방법이며, MPC 프로토콜이 검색의 실시간성에 끼치는 영향은 별도의 실험이 필요할 것이다.

\bigskip

결론적으로, 현재 상황에서 본 시스템이 지닌 보안 취약점을 완벽히 보완할 방법은 없어보인다. 어떤 방법을 선택하더라도 사용자 편의성의 희생, 검색의 실시간성의 손실 혹은 연산 비용의 증가는 필연적인 부작용으로 판단된다. 따라서 본 연구는 대규모 데이터 검색의 실시간성과 실용성을 달성하기 위해 OWNERSHIP 테이블 유지라는 전략적 타협을 선택하였다. 향후 위에서 제시한 더미 주입 및 비밀 분산 기법 등을 발전시킴으로써 부작용은 최소화 하며 최악의 보안 위협까지 완벽히 방어할 수 있는 범용적인 시스템으로 확장할 수 있기를 기대한다.